\chapter{СОСТАВ ЭКСПЕРИМЕНТАЛЬНЫХ АРТЕФАКТОВ}
\label{app:artefacts}
\sloppy

В таблице~\ref{tab:artefacts} приведён состав закоммиченных артефактов, на которых получены численные результаты главы~\ref{ch:results}. Артефакты сгруппированы по назначению и хранятся в каталоге \texttt{experiments/results/} репозитория проекта.

\begin{small}
\begin{longtable}{|>{\raggedright\arraybackslash}p{3.0cm}|>{\raggedright\arraybackslash}p{5.8cm}|>{\raggedright\arraybackslash}p{5.7cm}|}
\caption{Состав экспериментальных артефактов}\label{tab:artefacts}\\
\hline
\textbf{Группа} & \textbf{Назначение} & \textbf{Артефакты} \\
\hline
\endfirsthead
\caption*{Продолжение таблицы~\ref{tab:artefacts}}\\
\hline
\textbf{Группа} & \textbf{Назначение} & \textbf{Артефакты} \\
\hline
\endhead
\hline
\endfoot
\hline
\endlastfoot
Таблица единичных прогонов & Long-format таблица всех 486 единичных итераций параметрического симулятора с явным указанием значений всех параметрических осей конфигурации, идентификатора политики, номера реплики и значений показателей групп, описанных в подразделе~2.2.3 & \texttt{lhs\_parametric\_mobius\_2025\_autumn\_2026-05-08.\{json,csv,md\}} \\
\hline
Файлы постобработки параметрического прогона & Сводки попарных сравнений, агрегатов по уровням параметрических осей, диагностики risk-positive подмножества, чувствительности по осям, разброса между репликами, диагностики политики на основе языковой модели и распределения <<риск~$\times$~релевантность>> & \texttt{analysis\_pairwise.json}, \texttt{analysis\_capacity\_audit.json}, \texttt{analysis\_program\_effect.json}, \texttt{analysis\_gossip\_effect.json}, \texttt{analysis\_sensitivity.json}, \texttt{analysis\_stability.json}, \texttt{analysis\_llm\_ranker\_diagnostic.json}, \texttt{analysis\_risk\_utility.json} \\
\hline
Артефакты LLM-агентского симулятора & Long-format таблица 48 единичных итераций LLM-агентского симулятора на 12 точках репрезентативного подмножества & \texttt{llm\_agents\_lhs\_subset\_12pts.\{json,csv,md\}} \\
\hline
Артефакты перекрёстной проверки & Сводка коэффициентов ранговой корреляции Спирмена и совпадений верхнего лидера политик между параметрическим и LLM-агентским симуляторами на 12 общих точках, с пометками вырождения по точкам & \texttt{analysis\_cross\_validation.\{json,md\}} \\
\hline
Сводные отчёты & Компактные технические отчёты по результатам параметрического прогона (стадия~U) и совместного прогона с перекрёстной проверкой (стадия~W); содержат сводные таблицы и приёмочные чек-листы & \texttt{report\_mobius\_2025\_autumn\_parametric.md}, \texttt{report\_mobius\_2025\_autumn\_full.md} \\
\hline
Визуализации & Графические сводки результатов: диаграмма рассеяния <<риск~$\times$~релевантность>>, столбчатая диаграмма среднего риска по бакетам веса социального заражения, тепловая карта ранжирования политик на 12 maximin-точках, диаграмма коэффициентов ранговой корреляции по метрикам перекрёстной проверки & \texttt{plots/analysis\_risk\_utility\_scatter.png}, \texttt{plots/analysis\_gossip\_bucket\_bar.png}, \texttt{plots/analysis\_ranking\_heatmap\_maximin.png}, \texttt{plots/cross\_validation\_rho\_per\_metric.png} \\
\hline
Метаданные LLM-прогонов & Параметры запуска агентского симулятора, конкретные имена использованных языковых моделей (\texttt{openai/gpt-4o-mini} для политики на основе языковой модели, \texttt{openai/gpt-5.4-nano} для агентского симулятора), фактические значения суммарной стоимости и времени выполнения, число вызовов и доля ошибок парсинга структурированного ответа & Поля \texttt{params} и блок acceptance в \texttt{llm\_agents\_lhs\_subset\_12pts.json} \\
\hline
Тесты граничной верификации & Автоматизированные тесты, реализующие граничные свойства EC1--EC4 и их расширения на конфигурации с ненулевым весом социального заражения; каждый тест выполняется на минимально достаточной конфигурации с усреднением по нескольким случайным зёрнам & \texttt{experiments/tests/test\_extreme\_conditions.py} \\
\end{longtable}
\end{small}

Сохранённые артефакты позволяют воспроизвести любую единичную итерацию параметрического прогона по фиксированному значению параметрических осей и случайных зёрен. Воспроизводимость LLM-агентских прогонов ограничена стохастичностью внешней языковой модели и фиксируется на уровне условий запуска.
